\section{Threat Model}
\label{sec:threat}

\paragraph{System and attacker.}
Our target system is SonicBOOM (BOOMv3)~\cite{Zhao:SonicBOOM:2020}. It is an
open-source, out-of-order RV64GC core, and we use it inside a Chipyard 1.15.0
SoC~\cite{Amid:Chipyard:2020}. All programs, including the attacks, run bare-metal. We
simulate them cycle by cycle from the generated RTL in Verilator. We do not have a
physical FPGA board, so we use Vivado only to measure hardware cost
(Section~\ref{sec:approach}). We use the same attacker model as STT and
$\mu$STT~\cite{Yu:STT:2019,Chen:uSTT:2025}, which is the standard Spectre
setting~\cite{SpectreBase}. The attacker can watch any covert channel from anywhere a
thread can run. The attacker can also make the victim misspeculate, for example by
training the branch predictors or by choosing the victim's inputs. In our tests, the
attacker runs code in the same address space as the victim. The victim contains a
Spectre gadget, which is a load that reads a secret during misspeculation, followed by
an instruction whose execution depends on that secret. The attacker then recovers the
secret by timing accesses to a probe array in the L1 data cache with the
\texttt{rdcycle} counter. We use the BOOM Spectre proofs of concept for Spectre v1,
Spectre v2 and the return stack buffer variant.

\paragraph{Security goal.}
Our implementation is a speculation-aware defense. Following STT~\cite{Yu:STT:2019}, it
protects \emph{speculatively accessed} data, which is data returned by a load that may
later be squashed. It does not close any single covert channel. Instead, it stops any
transmitter that depends on such data from executing until the data is no longer
speculative (Section~\ref{sec:approach}). This means the guarantee covers the data
cache used by our attacks and every other channel a transmitter could use. We implement
the \emph{Spectre} model. In this model, only control-flow speculation is unsafe, and a
load stops being speculative once every older branch has resolved. The stricter
\emph{Futuristic} model moves the visibility point to the head of the reorder buffer,
and we treat it as a stretch goal. The defense succeeds if an attack that recovers the
secret on the baseline recovers nothing beyond chance on the protected configuration.

\paragraph{$\mu$STT under this threat model.}
$\mu$STT~\cite{Chen:uSTT:2025} keeps the threat model and security guarantee of STT and
only changes how it is built in hardware. It addresses two costs of STT, rename delay
and the area of the issue-queue delay logic, with the \emph{Age Matrix} and the
\emph{impede micro-op} (Section~\ref{sec:approach}). Neither mechanism has been built
into a complete core. The Age Matrix was synthesized only as a standalone Verilog
rename stage, where at a rename width of 8 it removes 85\% of the logic delay overhead
of STT and adds 36\% more area. The impede micro-op exists only in gem5, where $\mu$STT
has a 21\% overhead over an insecure baseline on SPEC CPU2017, compared with 16\% for
STT. Neither evaluation runs any attacks. Building both mechanisms into BOOM lets us
test the security guarantee on real RTL. It also raises a security question specific to
the impede micro-op. A transmitter that reaches execution without passing through the
issue queue, such as a load that is replayed or retried in the BOOM load/store unit,
would skip its impede micro-op. Our security evaluation must therefore cover every such
path.

\paragraph{Out of scope.}
As in STT, we do not protect data that the program reads non-speculatively. Leaking
such data is a traditional side channel, and constant-time programming addresses it.
Under the Spectre model, we also do not defend against transient execution that does
not come from branch misprediction, such as speculative store bypass or exception-based
attacks. Physical side channels, fault-injection attacks and attacks on the toolchain
are also out of scope. Our tests run bare-metal in a single address space, so the
victim and its secret model a software sandbox rather than a user/kernel privilege
boundary.
